\section*{第 \ref{ch:g5:heat-and-insulation} 章 --- 热与绝热}

\begin{solution}{exo:g5:heat-and-insulation:1}
热量是能量从较热的物体流向较冷的物体的过程；
它自己只会从热处流向冷处，
当两者温度相同时就停止流动。
\end{solution}

\begin{solution}{exo:g5:heat-and-insulation:2}
“关上门，你在让我们的热量跑出去！” --- 冷并不会流进来；
房子的热量会流向较冷的街道。
\end{solution}

\begin{solution}{exo:g5:heat-and-insulation:3}
快车道：铜锅、钢勺；
慢车道：羊毛围巾、静止的空气、木勺、羽绒被。
\end{solution}

\begin{solution}{exo:g5:heat-and-insulation:4}
锅底必须是快车道，把炉子的热量迅速传给食物；
锅柄必须是慢车道，把同样的热量挡在你的手外。
一件工具，两种本领，各得其所。
\end{solution}

\begin{solution}{exo:g5:heat-and-insulation:5}
两者温度相同，但金属是快车道：
它会迅速把你较暖的手中的热量带走，
这种快速散热的感觉就是“冰凉”。
皮肤感受的是热量散失的速度，而不是温度。
\end{solution}

\begin{solution}{exo:g5:heat-and-insulation:6}
公平的对比只改变一件事 --- 包裹物。
两只相同的杯子、相同的初始温度和同一个房间，
保证两条曲线的差异只能由围巾来解释。
\end{solution}

\begin{solution}{exo:g5:heat-and-insulation:7}
过了 $15$ 分钟：包着围巾的杯子约为 \qty{41}{\celsius}，
没包裹的杯子约为 \qty{30}{\celsius}。
两条曲线都趋向房间的 \qty{20}{\celsius}：
热量会一直流动到温度相同为止，保温只是让这段路程变慢。
\end{solution}

\begin{solution}{exo:g5:heat-and-insulation:8}
蓬松的羽毛在鸟的温暖身体周围
包裹着一层厚厚的静止空气 --- 那是最慢的慢车道。
同样的道理也用于羊毛衣服、羽绒被、
双层窗户（或抓绒衣、保温墙壁）。
\end{solution}

\begin{solution}{exo:g5:heat-and-insulation:9}
保温瓶既不制造热，也不制造冷，它只是封住热量的“车道”：
两层瓶壁之间是真空，内壁镀亮，把热辐射反射回去。
一月份，它减缓巧克力的热量向外流失；
七月份，它减缓夏天的热量向内渗入。
同一个器具两用，因为它能做的只是控制热量的流动。
\end{solution}

\begin{solution}{exo:g5:heat-and-insulation:10}
炎热的日子里，热量会从温暖的空气\emph{流入}
装冷饮的桶中。羊毛把这种向内的热流变成涓涓细流 ---
与它包裹温暖身体时的减缓作用相同。保温并不偏袒任何一方；
理解这一点的是卖家，而不是顾客。
\end{solution}

\begin{solution}{exo:g5:heat-and-insulation:11}
双层窗户在两层玻璃之间关住了一层静止空气 ---
一条慢车道，
于是房子的热量只能缓缓渗出，而不像单层窗户那样大量外流。
散失的热量少了，补充热量所需的燃料也就少了。
但保温只能减缓顺坡而下的热流，却永远无法完全阻止它：
冬天总会有一些热量漏出去，火炉必须不断地补充。
\end{solution}
